% Complete canonical interval and guarded-future proofs.

\section{The complete guarded future translation}
\label{sec:annex-p09-guarded-future-completeness}

\begin{annexitem}{annex:p9.gfc.label.transfer}{thm:p09-gfc-label-transfer}{The exact canonical-label substitution bridge}
\begin{proof}
Fix finitely many supplied formulas \(A_1,\ldots,A_s\)
of \(\mathcal A\). Choose fresh unary predicate letters
\(p_1,\ldots,p_s\), absent from those formulas. In a fixed
complete model interpret \(p_i\) by the entire truth set
of \(A_i\) in that model. Repeated occurrences of the same
input label may use the same placeholder; distinct placeholders
may also receive coincident truth sets.

The previously proved local theorems hold for every unary
interpretation, so they hold in this expanded model. Each
output label is obtained from the placeholders by finite
Boolean operations and strict future until. Replacing
\(p_i\) by \(A_i\) yields an actual \(\mathcal A\) formula.
Induction on temporal syntax shows that its truth is the
same as that of the placeholder formula in the expansion.
For first-order images, rename bound variables before
substituting each formula at its sole free current-point
argument. Induction on atoms, Boolean operations and
quantifiers gives the same substitution identity.

The infimum and test events are operations on the truth sets
of their labels. The substitution identity consequently
preserves their occurrence sets, definedness, values and
statuses, in order through their finite acyclic list.
Every infimum descriptor therefore retains its proved
meaning. The case compiler's transferred universal gap
conditions and point conditions retain their meanings too.
Its entire finite output becomes a canonical word list.

The positive-chain recursion uses no temporal constructor
outside Boolean operations and until on its input labels:
the neighborhood \(N^+A=A\mathbin{\mathsf U_s}\top\)
and the nonempty right limit
\(F\top\land\neg((\neg A)\mathbin{\mathsf U_s}\top)\)
have those explicit expansions. Its list-length induction
and finite-chain splice concern these truth sets in one
fixed model. They do not require the input label itself
to be future-dependent. Thus all cases of that already
proved construction transfer by the same substitution,
including a nonattained infimum equal to its lower endpoint.

Finally an output \(\alpha(a)\land D(a,b)\) in the former
left-labelled grammar becomes the present word with
first label \(\alpha\), the same interior labels and gaps,
and final label true. This is an exact syntactic inclusion.
No bound of the interval is inserted as a new free variable
inside a label. In particular a \(K^-\)-label may retain
information about the original model before its current
point; the model was not truncated in this argument.
\end{proof}
\end{annexitem}

\begin{annexitem}{annex:p9.gfc.positive.assembly}{thm:p09-gfc-positive-assembly}{Every merge, point label and concatenation condition}
\begin{proof}
Take two words at actual \(a<b\). Given realizations of
both, sort the union of their interior points. When a
point occurs in both chains, use one merged position
and conjoin its two labels. At the external endpoints
conjoin their endpoint labels as well.

For every original open gap, each point of that gap is
either one of the inserted merged points or belongs to
one of the smaller merged open gaps. Transfer its
original label to all those positions and subgaps,
excluding the two original endpoints. At each position
conjoin every imposed label. The original universal gap
condition is equivalent to precisely these transferred
conditions. This includes an empty original gap, which
contains neither an inserted point nor a subgap point.
Labels at a newly inserted point cannot be omitted.

Enumerate all order-preserving interleavings of the two
finite interior chains, allowing coincidence between
different chains. This is a finite list. The preceding
construction produces a word for each interleaving.
Any two original realizations supply one of its cases.
Conversely a merged realization restricted to either
original subsequence has all that word's point labels;
the transferred gap conditions recover every point of
each of its original gaps. It therefore realizes both
original words. This proves the conjunction identity
in both directions on every linear order.

For concatenation, suppose a left word ends at an actual
point \(c\) and a right word starts at that same point.
Join their point sequences, counting \(c\) once, and
conjoin both labels there together with any prescribed
extra midpoint label. All gaps on both sides retain their
original labels. This joined sequence gives a realization
on \((a,b)\). Conversely, splitting a joined realization
at its specified midpoint supplies the same \(c\) to
both constituent words. The strict inequalities make
that point interior. If the right side is a ray, its
last universal gap is simply retained.

Finite input disjunctions are handled by distributing all
pairs, and longer finite conjunctions by repeated merges.
The neutral word has no internal point and all three
labels true; the neutral ray has only its actual initial
point and a true final gap. Empty output lists mean false.
Every new label is a finite conjunction of supplied
\(\mathcal A\) formulas, so it remains in that class.
\end{proof}
\end{annexitem}

\begin{annexitem}{annex:p9.gfc.uniform.final.gap}{thm:p09-gfc-uniform-final-gap}{A completed handler without same-length negation recursion}
\begin{proof}
The event formula defining \(Q_W\) has exactly
\(n-1\) strictly internal required points. Its point
labels are \(A_1,\ldots,A_{n-1}\); its preceding
gap labels are \(B_1,\ldots,B_{n-1}\), in that
order. The fixed point \(a\) separately has
\(A_0\). Thus the transferred positive-chain
theorem identifies it exactly with a realization
of the indicated prefix ending at some \(z<b\).
It requires \(a<z\), even if all its labels are
true. An adjacent outer pair has no such chain.
Conjoin the \(A_0\) test at the lower root of
its event list and use the finite case compiler
on the negated Boolean decision. This gives the
explicit \(\overline Q_W\), including every
empty or nonattained event case. Hence this
handler has no unproved shorter-prefix interface.

Assume \(\mathsf{All}_{B_n}(a,b)\). For \(n\geq2\),
every realization of \(W\) restricts, by deleting
the last endpoint \(b\), to a realization of
its prefix ending at \(z_{n-1}<b\); it also
gives \(A_n(b)\). Conversely any such prefix
can be followed by \(b\): its new final gap
is contained in \((a,b)\) and therefore has
\(B_n\) throughout, and its last label is
the separately required \(A_n(b)\).
Consequently
\[
               [W](a,b)\longleftrightarrow A_n(b)\land Q_W(a,b)
\]
under the stated uniform-gap premise. Negating
gives exactly \(\mathsf{UC}_W\).

When \(n=1\), there is no internal point.
The word is \(A_0(a)\land A_1(b)\) under
that same gap premise, so its complement is
the separate one-interval clause. This is why
the zero-point prefix is not treated as an
existential interior witness. Each endpoint
failure is represented by a one-interval word
with true gap and the appropriate endpoint
label. The other alternatives are already
provided by the case compiler. All uses of
\(\mathsf{UC}\) retain the uniform-gap premise;
the handler alone is not claimed equivalent
to a general word's complement.
\end{proof}
\end{annexitem}

\begin{annexitem}{annex:p9.gfc.upper.boundary}{thm:p09-gfc-upper-boundary}{Every attained, unattained and terminal supremum}
\begin{proof}
Put \(C=\neg B_n\). The set
\(S=\{u:a<u<b,\ C(u)\}\) is empty
exactly when \(\mathsf{All}_{B_n}(a,b)\).
If nonempty it is bounded above by \(b\).
Completeness gives \(c=\sup S\), with
\(a<c\leq b\): any member of \(S\)
already lies strictly above \(a\).

We first verify the terminal case. If
\(c=b\), every \(r<b\) has a member
of \(S\) strictly above it; otherwise
\(r\) would be an upper bound smaller
than the supremum. This proves \(K^-C(b)\).
Conversely, apply \(K^-C(b)\) at the cut
\(a<b\) to obtain a member of \(S\).
For any \(r<b\), use the larger of
\(r,a\) as cut to obtain a member of
\(S\) above \(r\). Thus no point below
\(b\) bounds \(S\), and \(b\) is its
supremum. The predicate value at the
excluded endpoint \(b\) is irrelevant.

For \(a<c<b\), the supremum property
implies absence of \(C\) on \((c,b)\),
the displayed final-gap condition.
If \(C(c)\), then \(c\in S\) and
is its maximum. If \(\neg C(c)\),
every \(r<c\) again has an \(S\)-member
above it, giving \(K^-C(c)\).
These give the two interior subcases
and the disjunction in \(\Sigma_W\).

Conversely suppose \(\Sigma_W(a,c,b)\).
If \(C(c)\), this point belongs to
\(S\), and absence of later members
makes it the maximum. If instead
\(K^-C(c)\) supplies the disjunction,
its instance at \(a\) gives a member
of \(S\) below \(c\). Its instances
at cuts between \(a,c\) put members
arbitrarily close below \(c\).
The displayed absence beyond \(c\)
makes \(c\) an upper bound, and
those members show it is least.
If \(C(c)\) also happens to hold,
the attained subcase applies; if
not, the supremum is unattained.
Thus the two membership subcases
are disjoint, although the
descriptor's shorter disjunction
need not be.

An upper bound cannot be both
an interior point and \(b\).
Suprema in a linear order are
unique. The empty case is
disjoint from all nonempty
cases. These observations prove
the full classification. Every
use of \(K^-\) here is at
an actual point strictly above
\(a\); its instance at \(a\)
forces a strict interior
witness and removes any
first-point vacuity.
No assertion that an unattained
supremum is an occurrence
has been made.
\end{proof}
\end{annexitem}

\begin{annexitem}{annex:p9.gfc.word.cuts}{thm:p09-gfc-word-cuts}{Cutting and rejoining every finite realization}
\begin{proof}
Fix \(a<c<b\) and a realization
\(a=z_0<\cdots<z_n=b\) of \(W\).
Either \(c=z_i\) for a unique
internal index \(1\leq i<n\),
or \(z_{i-1}<c<z_i\) for a
unique gap index \(1\leq i\leq n\).
These are all positions because
the chain is finite and its
endpoints surround \(c\).

In the first case restrict
the realization at \(z_i=c\).
Its two restrictions realize
\(P_i\) and \(S_i\).
Both demand \(A_i(c)\);
this repetition is the
same point requirement,
not two separate witnesses.

In the second case insert
\(c\) into the original
gap. That gap supplies
\(B_i\) on both smaller
open gaps and at \(c\).
The resulting prefix is
exactly \(L_i\), and
the resulting suffix
is exactly \(R_i\).
This includes \(i=1\)
and \(i=n\). Their
extra endpoint label
\(B_i(c)\) is required
because \(c\) was an
interior point of
the original gap.

Conversely, a realizing
pair \(P_i,S_i\)
concatenates at \(c\)
to realize \(W\).
For a realizing pair
\(L_i,R_i\), delete
their shared point
\(c\) after concatenation.
The two adjoining
open \(B_i\)-gaps
and the \(B_i(c)\)
point label cover
the whole restored
gap. All remaining
point and gap
labels retain their
positions. The
result is a strict
realization of \(W\).
This proves both
directions of the
cut identity for
the chosen \(c\).
The proof did not
choose \(c\) as a
particular witness;
it therefore applies
to an independently
defined supremum.
\end{proof}
\end{annexitem}

\begin{annexitem}{annex:p9.gfc.interval.complement}{thm:p09-gfc-interval-complement}{The full mirrored negation construction and its strict descent}
\begin{proof}
We prove correctness and
finite construction together
by induction on the number
\(n\) of intervals. All
labels are supplied formulas
of \(\mathcal A\); their
internal complexity is
not an induction measure.
Assume actual \(a<b\).

At \(n=1\), the word
is exactly
\[
          A_0(a)\land A_1(b)\land
                         \mathsf{All}_{B_1}(a,b).
\]
Its negation is the
three alternatives in
the algorithm. The
gap failure has an
actual internal point
labelled \(\neg B_1\);
each side gap is true.
The two endpoint
failures need no
internal witness.
This proves the
base, including an
empty open interval.

For \(n\geq2\), the
three stated branch
types are exhaustive.
If \(A_n(b)\) fails,
\(E_1\) holds.
Otherwise either the
final-gap failure set
is empty, giving \(E_2\),
or it is nonempty.
If its supremum is
\(b\), the boundary
lemma gives the
other disjunct of
\(E_1\). If its
supremum is interior,
that lemma supplies
the unique \(c\)
and \(E_3(a,c,b)\).
The condition
\(\neg K^-(\neg B_n)(b)\)
in \(E_3\) retains
the distinction from
the terminal case.

Under \(E_1\), the word
is false. A failed
endpoint label plainly
prevents a realization.
If instead \(K^-(\neg B_n)(b)\)
holds, any proposed
last gap begins at
an actual \(z_{n-1}<b\).
The limit condition
at this cut supplies
a failure strictly
inside that gap,
contradicting its
universal \(B_n\).
The same reasoning
would also work
for a one-interval
word starting at \(a\).

Under \(E_2\), the
uniform-final-gap
lemma gives exactly
\(\mathsf{UC}_W\)
as the negation.
It has already been
constructed without
using this induction.
If \(a,b\) are
adjacent, \(E_2\)
or the endpoint
failure covers the
assignment; \(Q_W\)
has no realization,
so the returned
complement is true.
This handles empty
outer intervals
without inserting
an artificial point.

It remains to analyze
the interior-supremum
case. For its fixed
\(c\), the cut
identity gives
\[
 \neg[W](a,b)\ \longleftrightarrow\
 \bigwedge_{i=1}^{n-1}
       \bigl(\neg[P_i](a,c)\lor\neg[S_i](c,b)\bigr)
 \land
 \bigwedge_{i=1}^{n}
       \bigl(\neg[L_i](a,c)\lor\neg[R_i](c,b)\bigr).
 \tag{1}
\]
We justify every
replacement in this
finite conjunction.

For the point cuts,
both word lengths
are smaller than
\(n\). The induction
hypothesis therefore
replaces their
negations by
\(\mathsf N_{P_i}\)
and \(\mathsf N_{S_i}\).
Their endpoint
guards are supplied
by \(a<c<b\).
This gives \(D_i\)
for every index,
including \(i=1\)
and \(i=n-1\).

For gap cuts
\(2\leq i<n\),
the lengths
\(i\) and
\(n-i+1\)
are both less
than \(n\).
The induction
hypothesis gives
exactly \(H_i\).
When \(n=2\),
this range is
empty and its
conjunction is
true.

For the first
gap cut, \(L_1\)
has one interval
and is handled
by the base.
Its partner
\[
             R_1=(B_1,B_1,A_1,B_2,\ldots,B_n,A_n)
\]
has \(n\)
intervals. We
do not invoke
the induction
hypothesis on it.
The descriptor
\(\Sigma_W\)
supplies
\(\mathsf{All}_{B_n}(c,b)\).
Thus the completed
uniform handler
\(\mathsf{UC}_{R_1}(c,b)\)
is exactly
its negation
under the
retained descriptor.
This proves
the replacement
by \(H_1\)
with no
same-length
recursion.

For the last
gap cut, \(L_n\)
has \(n\)
intervals but
is impossible
under \(\Sigma_W\).
It requires
\(B_n(c)\).
If \(\neg B_n(c)\)
holds in the
descriptor, that
point requirement
already fails.
Otherwise the
descriptor has
\(K^-(\neg B_n)(c)\).
A realization
of \(L_n\)
would have
a last gap
ending at
\(c\), beginning
at a strictly
earlier actual
point, with
\(B_n\)
throughout.
The past limit
supplies a
failure in
that very
gap. Hence
\(L_n\)
is false
in this
case too.
Its entire
factor
\(\neg L_n\lor\neg R_n\)
in (1) is
therefore true
and may be
omitted.
This proves
the algorithm's
last extreme
replacement,
retaining the
actual descriptor
that licenses it.

The remaining
conjunction in
(1) is exactly
\(M_W(a,c,b)\).
Thus, for this
unique interior
boundary, it
expresses the
word's negation.
Conversely any
witness for
\(E_3\land M_W\)
satisfies the
same descriptor,
so all the
reverse implications
just proved
apply and
give \(\neg W\).
Existentially
quantifying \(c\)
therefore gives
precisely the
third branch
of the output.
Together with
the first two
branches this
proves the
required
equivalence.

We finish by
making the
compilation
claim explicit.
Distribute the
finite alternatives
in every \(D_i\)
and \(H_i\)
and the
finite lists
already supplied
by recursive
calls and
uniform handlers.
Each chosen
term is a
conjunction
of point
conditions at
\(a,c,b\)
and word
conditions on
their intervals.
A word on
\((a,b)\)
is split
at \(c\)
using the
proved finite
cut identity
and all
its alternatives.
For a
one-interval
word there
is the
single gap
cut: its
gap label
must also
hold at
\(c\).
Conjoin
the resulting
lists on
\((a,c)\)
by the
positive merge
construction,
and do the
same on
\((c,b)\).
Conjoin all
point conditions
into the
appropriate
endpoint labels.
Concatenate
at \(c\).
The merged
realization
equivalences
prove this
compilation
in both
directions.
Thus no
untranslated
three-variable
condition
remains.

Every full
\(\mathsf N\)
call made
at length
\(n\) has
length at
most \(n-1\).
The uniform
handlers use
their independently
terminating
positive-chain
event recursion.
All conjunctions,
cut lists
and merges
are finite.
Induction
therefore proves
termination of
the entire
procedure.
The only
new past
operator is
the explicitly
allowed \(K^-\)
at the
supremum.
All other
added label
operations
are Boolean
or strict
until.

If the
endpoint
values are
not ordered,
the word
is false.
The unguarded
complement
therefore
adds equality
and reversed
order as
two separate
cases. On
the ordered
domain, the
complement of
a finite
disjunction
is the
finite
conjunction
of these
constructed
complements,
and positive
assembly
compiles it.
This proves
the full
canonical
relative
closure.
\end{proof}
\end{annexitem}

\begin{annexitem}{annex:p9.gfc.ray.translation}{thm:p09-gfc-ray-translation}{Every finite ray constructor and its terminal cases}
\begin{proof}
At its final chosen point \(z_m\), the condition is
\(A_m(z_m)\) and the truth of \(B_{m+1}\) at every
strictly later point. This is exactly \(T_m\).
Inducting backwards, a realization beginning at
\(z_i\) requires \(A_i(z_i)\), one strict future
point \(z_{i+1}\) realizing its suffix, and \(B_{i+1}\)
on the open gap. These are precisely the clauses of
\(T_i\). Conversely its strict until witness supplies
the next point and the induction hypothesis supplies
the remainder of the ray realization. This proves
the displayed translation in both directions.

All labels retain their whole-model meaning. A past
limit inside a label is evaluated at that actual point,
not at a formal endpoint or in an independently
truncated submodel. The backward recursion has
finitely many steps. Its substitution of each supplied
label therefore gives a finite \(\mathcal A\) formula.

For a cofinal conjunct
\(\forall u\in T\,\exists v>u\,\Psi(v)\), the formula
\(F\Psi\land GF\Psi\) at \(x\) is exact. The
first factor supplies a witness above \(x\), which
also works for every \(u\leq x\). For \(u>x\),
the second factor supplies its later witness.
Conversely the cofinal assertion supplies the
first factor by taking \(u=x\) and the second
by taking each \(u>x\). This retains the
nonempty factor even at a last point.

These formulas cover all terminal possibilities.
If there is no failure of the final gap after
the current point, its \(G\) condition is
automatic at every later last witness.
If the failures are cofinal, then every
proposed final witness has a later failure,
so its \(G\) condition and the ray word
are false. Cofinality here includes an
unbounded future failure set; no supremum
point is introduced for it. If the future
failure set is nonempty and bounded above
in a complete order, let \(c\) be its
supremum. A possible final witness below
\(c\) has a failure after it and cannot
satisfy \(G\). A witness at or above
\(c\) has no later failure and satisfies
that terminal gap, whether \(c\) is a
failure itself or an unattained limit.
The point label still has to hold at
the chosen witness. This describes the
attained and unattained cases without
requiring a label at infinity.

At a last point, a zero-witness ray has
value \(A_0(x)\), because its only gap
universal is vacuous. A positive-witness
ray is false, because its first strict
until has no witness. A required cofinal
conjunct is false there, due to \(F\Psi\).
The same values follow from the original
first-order ray conditions.

Finite lists translate by disjunction.
If \(B(x)\) is their resulting formula,
the zero-witness ray with point label
\(\neg B\) and true final gap has value
exactly \(\neg B(x)\). Empty input
lists and empty conjunctions are thus
handled without any further construction.
\end{proof}
\end{annexitem}

\begin{annexitem}{annex:p9.gfc.regional.induction}{thm:p09-gfc-regional-induction}{The signed regional induction and exact normal-form bridge}
\begin{proof}
Use the literal GPSS regional stages from
\cref{def:p09-gpss-regional-stages}.
Their level zero is empty. At a subsequent
level a component is either a guarded
existential with finitely many children,
or its negation. Each child is an earlier
left or right regional component, or
one signed unary literal at the new
interior point. In the applications here
the left endpoint is always actual;
the right endpoint is actual or the
formal future endpoint.

Assume by induction that lists are
available for every earlier regional
child. At a positive constructor,
distribute its finite conjunction over
all choices of child disjuncts.
Use positive assembly for the chosen
left children on the left interval
and the chosen right children on
the right interval. All point
literals conjoin at the quantified
point. Concatenate the two lists
there. If either family is empty,
use its true neutral interval or
ray; if both are empty, the
strict existential point remains
required. If any child has an
empty list, that distributed
choice has no realization and
contributes no output.

This construction works equally
for a bounded outer interval
and a future ray. The midpoint
is actual, so its literals and
the two endpoint labels from
concatenation all concern the
same domain point. Endpoint
coincidence is already excluded
by the constructor's strict
guard. No additional point
is inserted into an empty
open interval.

For a negative constructor,
first construct the list for
its underlying positive
existential. On a bounded
interval, complement every
word by the full canonical
complement theorem and
conjoin the resulting lists.
On a ray, apply the complete
ray complement translation.
An empty positive list is
false and its complement
is the neutral true list.
These steps handle both
signs of every source
constructor. They use
only complete-order
complement equivalences,
under the actual strict
interval guards. The
stage induction is
therefore valid at
every finite level.

For the final first-order
bridge, take an input
whose quantifiers are
restricted to the
current tail. The
already proved GPSS
future special normal
form yields a finite
disjunction of conjunctions
of \(\operatorname{At}_G(x)\)
and \(\operatorname{Aft}_G(x)\)
components. Its signed
point literals are
atomic or negated
atomic \(\mathcal A\)
labels. The induction
just completed supplies
canonical ray lists
for every after
component; positive
assembly combines
each finite clause.
The ray translation
then gives an actual
\(\mathcal A\) formula.

This bridge relates
the different source
grammars explicitly.
GPSS uses only
signed atomic
point children
in its initial
regional normal
form. Rabinovich's
canonical word
labels may be
arbitrary supplied
\(\mathcal A\)
formulas, including
at both actual
endpoints. The
former point
grammar embeds
in the latter;
the latter labels
are generated
by the proved
complement
construction,
not imported
as arbitrary
one-variable
first-order
conditions.
All labels
are inlined
by finite
substitution.

The GPSS normal
form already
includes every
weak order type
and every quantified
point position,
including equality
with the current
point. Its use
here therefore
provides a complete
assembly in place
of the source's
abbreviated final
structural induction.
It does not require
before components
for the future
input, and it does
not infer unrestricted
canonical expressibility
of unbounded-from-below
sentences. The
canonical interval
and ray results
used are exactly
the ones required
by this stage
construction.
\end{proof}
\end{annexitem}

\begin{annexitem}{annex:p9.guarded.future.completeness}{thm:p09-guarded-future-completeness}{The effective translation, initial normalization and converse}
\begin{proof}
Let \(\varphi(x)\) be
a future monadic
first-order formula.
Rename all its bound
variables away from
\(x\). Restrict
each quantifier to
the tail \([x,+\infty)\).
By induction on
first-order syntax,
the resulting formula
\(\varphi^+(x)\)
has in the original
model the value of
\(\varphi\) in that
induced tail.
Because \(\varphi\)
is future, these
values agree at
the current point.

The displayed
inclusive-to-strict
identities follow
by dividing a
quantified value
\(y\geq x\)
into its equality
case \(y=x\)
and strict case
\(y>x\).
For an existential,
either case may
provide its
witness. For
a universal,
both cases
are required.
Fresh bound
renaming makes
the substitution
\(\chi[x/y]\)
capture-free.
Applying these
finite identities
recursively gives
the strict source
presentation when
desired, while
the GPSS normal
form can use
the inclusive
version directly.
Thus no present
atomic condition
is discarded.

Apply the complete
regional construction
to obtain
\(\theta\in\mathcal A\)
equivalent to
\(\varphi\)
over complete
orders. Each
step is finite:
normal-form
distribution,
the regional
stage construction,
strictly decreasing
word-complement
recursion, terminating
positive-chain
event construction,
finite merges,
and finite
label substitution.
No semantic
search over
models is
part of the
procedure.

We next verify
the final
normalization.
At a first
point \(x\),
the displayed
universal
definition of
\(K^-B(x)\)
is true for
every \(B\),
since there
is no earlier
cut. Induction
on the outer
Boolean syntax
therefore gives
\[
                  \theta(x)\longleftrightarrow g(\theta)(x)
\]
at every first
point. An
until node
is retained
as a whole;
no claim
is made that
its operands
are evaluated
at that first
point. Every
surviving
past-limit
occurrence
is beneath
such an
until node,
so \(g(\theta)\)
is in the
guarded future
fragment.

For any
complete model
\(\mathcal M\)
and point
\(x\), the
tail
\(\mathcal N=\mathcal M|_{\geq x}\)
is complete.
Indeed a
nonempty subset
of that tail
has its
infimum in
the tail
whenever
bounded below:
\(x\) is
also a lower
bound in
\(\mathcal M\),
so its
infimum is
at least
\(x\).
If it is
bounded above,
its supremum
is at least
one of its
members, hence
at least \(x\).
The original
complete-order
bounds therefore
belong to
the induced
tail and have
the same
extremal
properties
there.

Now \(\varphi\)
is future,
and both
\(\mathcal M\)
and \(\mathcal N\)
are in the
class where
the constructed
\(\theta\)
is equivalent
to it.
The guarded
future lemma
\cref{thm:p09-hr-guarded-future}
applies to
\(g(\theta)\)
on all
linear orders.
Consequently
\[
\begin{aligned}
 \mathcal M,x\models\varphi
 &\Longleftrightarrow \mathcal N,x\models\varphi\\
 &\Longleftrightarrow \mathcal N,x\models\theta\\
 &\Longleftrightarrow \mathcal N,x\models g(\theta)\\
 &\Longleftrightarrow \mathcal M,x\models g(\theta).
\end{aligned}
\]
The third
step uses
that \(x\)
is first
in \(\mathcal N\);
the last
uses the
guarded
fragment's
proved
future
semantics.
This proves
the forward
translation
at every
actual point,
including a
first point,
a last point
or the point
of a singleton
order.

For the converse,
expand a
guarded temporal
formula into
first-order
syntax. At
a named
current point
\(v\), the
two non-Boolean
clauses are
\[
\begin{aligned}
 (A\mathbin{\mathsf U_s}B)^*(v)
   &=\exists w\bigl(v<w\land B^*(w)\land
             \forall u(v<u<w\to A^*(u))\bigr),\\
 (K^-A)^*(v)
   &=\forall r\bigl(r<v\to
             \exists u(r<u<v\land A^*(u))\bigr).
\end{aligned}
\]
Use fresh
bound variables
at every
copied node;
atomic and
Boolean clauses
commute with
translation.
Induction on
the temporal
tree proves
its correctness
on every
linear-order
model. By
the guarded
future lemma
this image
is semantically
future. Restricting
all its
quantifiers
to the
current tail
therefore
preserves its
truth and
gives the
explicit future
first-order
formula.
The converse
needs no
completeness
assumption.

For the
discrete
specialization,
every point
strictly after
the current
one has
an immediate
predecessor.
Its open
interval from
that predecessor
is empty,
so \(K^-B\)
is false
there regardless
of \(B\).
Traverse the
guarded output
through its
Boolean and
until nodes,
stopping whenever
a past-limit
node is reached.
Each such node
has only those
future and Boolean
ancestors, with
at least one
strict until,
so it is
evaluated strictly
after the current
point. Replace
that entire
subformula by
false without
traversing its
operand.
Induction over
the remaining
formula gives
an equivalent
\(U\)-only
formula on
discrete
complete
orders.
No replacement
of an
unguarded
past limit
at a first
point is
being made.

Finally consider
the body's
attainment
examples. With
a single
interior witness
at 1 and
final failures
only at 1,
the final
open gap
\((1,2)\)
has no
failure.
The point
cut at 1
therefore
retains a
valid realization,
despite that
point being
the attained
failure supremum.
With failures
on \((1/2,1)\)
the same
open final
gap has
none, while
the supremum
1 is
unattained.
The descriptor
now uses
the past-limit
disjunct.
Failures
cofinal below
2 contradict
every possible
last gap
by its
instance at
that last
witness.

For \(F(P\land K^-P)\),
the eventuality
witness is
strictly after
the current
point. The
past-limit
condition there
can be checked
using only
cuts and
witnesses between
the current
point and
that later
point, by
the same
guarded-future
argument.
Conjoining \(P\)
at the current
point yields
the displayed
future limit
counterexample
to \(U\)-only
expressiveness.
The positive
basis expresses
it precisely
because the
required
past-limit
operation is
permitted
inside the
strict future
guard.
\end{proof}
\end{annexitem}
